% Partie : intuitionistic-logic
% Chapitre : introduction
% Section : bhk-interpretation

\documentclass[../../../include/open-logic-section]{subfiles}

\begin{document}

\olfileid[fr]{int}{int}{bhk}

\olsection{L'interprétation de Brouwer--Heyting--Kolmogorov}

Les constructions des exemples suivent la forme des formules.
Un couple fournit les deux constructions d'une conjonction.
Une disjonction exige le choix d'un membre et une construction
de ce membre. Une implication demande un procédé qui transforme
une construction de l'antécédent en une construction du conséquent.

Les connecteurs intuitionnistes admettent une interprétation
constructive informelle, généralement appelée interprétation de
Brouwer--Heyting--Kolmogorov. Elle repose sur la notion de
« construction », que l'on peut comprendre comme une preuve
constructive. Nous évitons toutefois le mot « preuve » dans cette
interprétation pour ne pas confondre cette notion avec celle d'%
!!a{derivation} dans un système formel de !!{derivation}.
À partir de cette notion intuitive, l'interprétation BHK explique
le sens des connecteurs intuitionnistes. Les fonctions utilisées
ci-dessous sont des procédés effectifs de transformation des
constructions ; il ne suffit pas d'affirmer l'existence d'une
fonction ensembliste quelconque.

\begin{enumerate}
\item Nous supposons savoir ce qui constitue une construction
  d'un énoncé atomique.
\item Une construction de $!A_1 \land !A_2$ est un couple
  $\tuple{M_1, M_2}$, où $M_1$ est une construction de~$!A_1$
  et $M_2$ une construction de~$!A_2$.
\item Une construction de $!A_1 \lor !A_2$ est un couple
  $\tuple{s, M}$, où soit $s$ vaut~$1$ et $M$ est une construction
  de~$!A_1$, soit $s$ vaut~$2$ et $M$ est une construction de~$!A_2$.
\item Une construction de $!A \lif !B$ est une fonction qui
  transforme une construction de~$!A$ en une construction de~$!B$.
\item Il n'existe aucune construction de~$\lfalse$ (l'absurde).
\item $\lnot !A$ est défini comme synonyme de $!A \lif \lfalse$.
  Autrement dit, une construction de $\lnot !A$ est une fonction
  qui transforme une construction de~$!A$ en une construction
  de~$\lfalse$.
\end{enumerate}

\begin{ex}
Prenons $\lnot \lfalse$. Une construction de cette formule est
une fonction qui, recevant une construction de $\lfalse$, renvoie
une construction de~$\lfalse$. La fonction identité~$\Id{}$
convient : si l'on lui donne une construction~$M$ de~$\lfalse$,
alors $\Id{}(M) = M$ fournit une construction de~$\lfalse$.
Cette condition sur les entrées n'affirme pas qu'une telle entrée
existe : l'absurde n'a aucune construction.
\end{ex}

De manière générale, $\lnot !A$ signifie : « Une construction
de $!A$ est impossible ». Cette impossibilité est établie par un
procédé conduisant à l'absurde, et ne se réduit pas à l'absence
d'une construction actuellement connue.

\begin{ex}
Donnons une construction de $!A \lif \lnot \lnot !A$, pour toute
proposition~$!A$, c'est-à-dire de
$!A \lif ((!A \lif \lfalse) \lif \lfalse)$.
Il faut une fonction~$f$ qui, recevant une construction~$M$
de $!A$, renvoie une construction~$f(M)$ de
$(!A \lif \lfalse) \lif \lfalse$.
Pour définir $f(M)$, qui doit être une construction de
$(!A \lif \lfalse) \lif \lfalse$, nous définissons
une fonction $g$ qui, recevant une construction~$h$
de $!A \lif \lfalse$, renvoie une construction de~$\lfalse$.
La fonction $g$ applique~$h$ à la construction~$M$ de $!A$
reçue auparavant. Le résultat~$h(M)$ de $h$ est alors une
construction de $\lfalse$. Ainsi, $f(M)(h) = h(M)$ est une
construction de~$\lfalse$ dès lors que $M$ est une construction
de~$!A$. La fonction renvoyée par la première étape attend donc
la construction d'une négation comme deuxième entrée.
\end{ex}

\begin{ex}
Donnons une construction de $\lnot(!A \land \lnot !A)$,
c'est-à-dire de $(!A \land (!A \lif \lfalse)) \lif \lfalse$.
Il s'agit d'une fonction~$f$ qui, recevant une construction~$M$
de $!A \land (!A \lif \lfalse)$, renvoie une construction
de~$\lfalse$. Une construction d'une conjonction
$!B_1 \land !B_2$ est un couple $\tuple{N_1, N_2}$, où
$N_1$ est une construction de~$!B_1$ et $N_2$ une construction
de~$!B_2$. Les fonctions $p_1$ et $p_2$ extraient d'une
construction de $!B_1 \land !B_2$ les constructions de
$!B_1$ et de $!B_2$, respectivement :
\begin{align*}
  p_1(\tuple{N_1, N_2}) &= N_1\\
  p_2(\tuple{N_1, N_2}) &= N_2
\end{align*}
Voici ce que fait $f$. Il applique d'abord $p_1$ à son entrée~$M$,
ce qui donne une construction de~$!A$. Puis il applique $p_2$
à~$M$, ce qui donne une construction de~$!A \lif \lfalse$.
Cette dernière est une fonction~$p_2(M)$ qui, recevant une
construction de~$!A$, renvoie une construction de~$\lfalse$.
L'application de $p_2(M)$ à $p_1(M)$ donne donc une construction
de~$\lfalse$. Nous pouvons ainsi définir
$f(M) = p_2(M)(p_1(M))$.
\end{ex}

\begin{ex}
Donnons une construction de $((!A \land !B) \lif !C) \lif (!A \lif
(!B \lif !C))$ : une fonction $f$ transformant une construction~$g$
de $(!A \land !B) \lif !C$ en une construction de
$(!A \lif (!B \lif !C))$. La construction $g$ est elle-même
une fonction, qui transforme des constructions de $!A \land !B$
en constructions de~$!C$. Son image $f(g)$ doit être une
fonction~$h_g$ qui transforme des constructions de $!A$ en
fonctions allant des constructions de~$!B$ aux constructions de~$!C$.

Décomposons cette construction en trois étapes. Nous devons définir
une fonction $h_g$, qui sera le résultat de $f$ appliqué à~$g$.
L'entrée de~$h_g$ est une construction~$M$ de~$!A$.
Le résultat de $h_g(M)$ doit être une fonction~$k_{g,M}$
transformant des constructions~$N$ de~$!B$ en constructions
de~$!C$. Posons $k_{g,M}(N) = g(\tuple{M, N})$.
Le couple $\tuple{M, N}$ est une construction de~$!A \land !B$.
Ainsi, $k_{g,M}$ est une construction de $!B \lif !C$ : elle
transforme les constructions~$N$ de~$!B$ en constructions de~$!C$.
Posons ensuite $h_g(M) = k_{g,M}$. Cette fonction associe aux
constructions~$M$ de~$!A$ des constructions $k_{g,M}$ de
$!B \lif !C$. Posons enfin $f(g) = h_g$. Cette fonction
transforme les constructions $g$ de $(!A \land !B) \lif !C$
en constructions de $!A \lif (!B \lif !C)$.
Les trois paramètres sont introduits successivement ; les fonctions
intermédiaires conservent les paramètres déjà reçus.
\end{ex}

L'énoncé $!A \lor \lnot !A$ exprime le principe du tiers exclu.
On peut en construire certaines instances, pour des propositions
$!A$ particulières, par exemple $\lfalse \lor \lnot \lfalse$,
mais on ne dispose pas d'une construction intuitionniste uniforme
de toutes ses instances. La construction d'une disjonction doit
indiquer quel membre est construit ; on ne peut donc obtenir ce
principe simplement en affirmant que toute proposition est vraie
ou fausse. Les énoncés pour lesquels nous ne connaissons ni
démonstration ni réfutation, comme la conjecture binaire de Goldbach,
illustrent cette différence, sans constituer à eux seuls une preuve
de non-dérivabilité. En revanche, la double négation du tiers exclu
est constructible : $\lnot \lnot (!A \lor \lnot !A)$ est valide.

\begin{ex}
Pour construire $\lnot \lnot (!A \lor \lnot !A)$, il nous faut
une fonction~$f$ qui transforme une construction de
$\lnot(!A \lor \lnot !A)$, c'est-à-dire de
$(!A \lor (!A \lif \lfalse)) \lif \lfalse$, en une construction
de~$\lfalse$. Autrement dit, il faut une fonction $f$ telle que
$f(g)$ soit une construction de~$\lfalse$ dès lors que $g$ est
une construction de $\lnot(!A \lor \lnot !A)$.

Supposons que $g$ soit une construction de $\lnot(!A \lor \lnot !A)$,
donc une fonction transformant une construction de $!A \lor \lnot !A$
en une construction de~$\lfalse$. Une construction de $!A \lor \lnot !A$
est un couple $\tuple{s, M}$, où soit $s = 1$ et $M$ est une
construction de $!A$, soit $s = 2$ et $M$ est une construction
de $\lnot !A$. Soit $h_1$ la fonction transformant une construction~$M_1$
de $!A$ en une construction de $!A \lor \lnot !A$ : elle
associe à $M_1$ le couple $\tuple{1, M_1}$.
Soit $h_2$ la fonction transformant une construction~$M_2$
de $\lnot !A$ en une construction de $!A \lor \lnot !A$ :
elle associe à $M_2$ le couple $\tuple{2, M_2}$.

Posons $k = \comp{h_1}{g}$. Cette fonction transforme une
construction de~$!A$ en une construction de~$\lfalse$ : c'est
donc une construction de~$!A \lif \lfalse$, ou de $\lnot !A$.
Posons ensuite $l = \comp{h_2}{g}$. Cette fonction transforme
une construction de $\lnot !A$ en une construction de~$\lfalse$.
Comme $k$ est une construction de $\lnot !A$, le résultat
$l(k)$ est une construction de~$\lfalse$.

Nous avons ainsi expliqué comment transformer une construction~$g$
de $\lnot(!A \lor \lnot !A)$ en une construction de~$\lfalse$.
La fonction $f$ associant à une construction~$g$ de
$\lnot(!A \lor \lnot !A)$ la construction $l(k)$ de~$\lfalse$
est donc une construction de $\lnot\lnot(!A \lor \lnot !A)$.
\end{ex}

Comme on le voit, utiliser l'interprétation BHK pour établir la
validité intuitionniste de !!{formula}s devient vite laborieux
et difficile à suivre. Heureusement, il existe des systèmes de
!!{derivation} plus commodes pour la logique intuitionniste,
ainsi que des interprétations sémantiques plus précises.

\end{document}
